\begin{abstract}
Sample-efficient pretraining under a developmental data budget, such as
the ten million words and ten epochs of the BabyLM Strict-Small track, is
usually approached by adding training signal: a new objective, optimizer,
curriculum, or teacher. The budget itself receives less scrutiny. We show
that much of it is spent after learning has stopped. A
\rParamsM{}-parameter GPT-BERT hybrid trained under the track rules
saturates on downstream evaluations by roughly \rPeakCkpt{} words seen,
about $70\%$ of its permitted exposure, while training loss continues to
fall; the final epochs improve fit, not competence. From this diagnosis we
derive three design principles, timing, preservation, and alignment, and
instantiate them as \emph{\MPOlong{}} (\MPO{}): a short post-training
phase, begun at the saturation checkpoint, that teaches the model to
prefer real corpus sentences over length-preserving corruptions of them
under a frozen-reference penalty. A controlled study of seven in-training
interventions supports the principles: none improves on the baseline, and
the unanchored ones damage the same cluster of long-distance dependencies
first. \MPO{} improves five of the seven official task averages, concedes
under a point on the other two, is alone in repairing that cluster, and
matches the fully trained baseline with \rDpoEpochs{} rather than
\rEpochs{} epochs; our \MPO{} entry places \rRankDpo{} in the track.
\end{abstract}
